\section{Discretization of Gaussians}

We investigate a discretization scheme for Gaussians, which preserves the maximum entropy property with respect to the sufficient statistic.

\subsection{Standard Gaussians}

A Gaussian distribution on $\rr$ with mean $\seccatindex$ and standard derivation $\sigma$ is the distribution
\begin{align*}
    \probat{\indexedcatvariable}
    &= \frac{1}{\sqrt{2\pi\sigma^2}} \expof{-\frac{(\catindex-\seccatindex)^2}{2\sigma^2}} \\
    &= \expof{- \frac{1}{\sigma^2} \catindex^2 + \frac{\seccatindex}{\sigma^2} \catindex - \frac{\seccatindex^2}{\sigma^2} - \lnof{\sqrt{2\pi\sigma^2}}} \, .
\end{align*}

We notice that it is in the exponential family to the statistic $\sstat(\catindex)=(\catindex,\catindex^2)$ with the canonical parameter
\begin{align*}
    \canparamat{\selvariable}
    = \coloredmatrixof{
    \frac{\seccatindex}{\sigma^2} \\
    -\frac{1}{\sigma^2}
    }{\selvariable}
\end{align*}
and the mean parameter (can be verified by integration)
\begin{align*}
    \meanparamat{\selvariable}
    = \coloredmatrixof{
    \seccatindex \\
    \sigma^2+\seccatindex^2
    }{\selvariable} \, .
\end{align*}

The mean polytope is the area bounded by a parabola (the image set of the statistic).
The zero-noise limit are the maximum entropy distributions on the parabola, namely the Dirac distributions to the corresponding values.

\subsection{Discretization}

Let us now restrict the values $\rr$ to a finite set $\arbset\subset\rr$, which is enumerated by a variable $\catvariable$ taking values in $[\catdim]$.
We define discrete Gaussian distributions on $\arbset$ as members of the exponential family with sufficient statistic $\sstat(\catindex)=(\catindex,\catindex^2)$.
The mean polytope is now the convex hull of the two-dimensional vectors on the parabola which are images of $\arbset$.


\subsection{Multivariate Gaussians}

The discretization scheme preserving the maximum entropy property can be extended to multivariate Gaussians.
The sufficient statistics is then $\sstat(\shortcatindices)=(\shortcatindices,\shortcatindices\shortcatindices^T)$, where by $\shortcatindices\shortcatindices^T$ we denote the unfolded matrix of product of the indices.